% Supplementary material of the Information Sciences submission (doc/infsci/main.tex).
% Every table is generated by scripts/_emit_infsci_tables.py into generated/, except Table S2, which lists design choices.
% Cross-references to the main text resolve through xr: build with scripts/_build_infsci.sh.
\documentclass[preprint,12pt]{elsarticle}
\usepackage[T1]{fontenc}
\usepackage[utf8]{inputenc}
\usepackage{lmodern}
\usepackage{amsmath}
\usepackage{amssymb}
\usepackage{booktabs}
\usepackage{array}
\usepackage{microtype}
\usepackage{xurl}
\usepackage{graphicx}
\usepackage{float}
\usepackage{xr}
\usepackage[section]{placeins}
\externaldocument[M-]{main}
\renewcommand{\thetable}{S\arabic{table}}
\renewcommand{\thesection}{S\arabic{section}}
\renewcommand{\theequation}{S\arabic{equation}}
\renewcommand{\topfraction}{0.9}
\renewcommand{\bottomfraction}{0.8}
\renewcommand{\textfraction}{0.07}
\renewcommand{\floatpagefraction}{0.8}
\setcounter{topnumber}{3}
\setcounter{totalnumber}{4}
\newcolumntype{L}[1]{>{\raggedright\arraybackslash}p{#1}}
\newcommand{\partheading}[1]{\bigskip\noindent{\large\bfseries #1}\par\nobreak\smallskip}
% generated by scripts/_emit_infsci_tables.py; do not edit
\newcommand{\floatMaxRP}{1.4e-13}
\newcommand{\floatMaxGASF}{5.6e-17}
\newcommand{\floatMaxGADF}{8.4e-17}
\newcommand{\floatMaxXIRP}{2.8e-16}
\newcommand{\floatBlocksExactRP}{19}
\newcommand{\rpEightMin}{0.0009}
\newcommand{\rpEightMax}{0.0022}
\newcommand{\rpEightSwarmMin}{0.012}
\newcommand{\rpEightSwarmMax}{0.017}
\newcommand{\benchSeries}{54\,336}
\newcommand{\benchBlocks}{20}
\newcommand{\macroDiffRPHIRATA}{0.185}
\newcommand{\macroDiffRPHIRATASIGNED}{0.084}
\newcommand{\macroDiffRPUNSIGNED}{0.122}
\newcommand{\macroDiffMTF}{0.311}
\newcommand{\binaryOracleMin}{0.03}
\newcommand{\binaryOracleMax}{0.17}
\newcommand{\condGoodSeries}{47\,949}
\newcommand{\condGoodBad}{0}
\newcommand{\condWorstRate}{85}
\newcommand{\speedupMid}{2.9}
\newcommand{\speedupLong}{5.2}
\newcommand{\clsCombos}{84}
\newcommand{\clsCombosEq}{66}
\newcommand{\clsWindows}{90\,937}
\newcommand{\escWindows}{18\,000}
\newcommand{\escCombos}{36}
\newcommand{\robHighRP}{0.041}
\newcommand{\robHighGASF}{0.064}
\newcommand{\robHighGADF}{0.064}
\newcommand{\robHighXIRP}{0.011}
\newcommand{\robDatasetsRPBetter}{6}
\newcommand{\robDatasets}{7}
\newcommand{\robLossyShift}{0.019}
\newcommand{\srRPFivehundred}{0.0076}
\newcommand{\srColFivehundred}{0.08}
\newcommand{\srShareFivehundred}{67}
\newcommand{\srRPTwofifty}{0.0041}
\newcommand{\srColTwofifty}{0.16}
\newcommand{\srShareTwofifty}{23}
\newcommand{\srRPHundred}{0.0015}
\newcommand{\srColHundred}{0.37}
\newcommand{\srShareHundred}{0}
\newcommand{\srRPFifty}{0.0009}
\newcommand{\srColFifty}{0.61}
\newcommand{\srShareFifty}{0}
\newcommand{\srGafMin}{0.0018}
\newcommand{\srGafMax}{0.0032}
\newcommand{\srNativeRP}{0.0015}
\newcommand{\srResampledRP}{0.0015}


\begin{document}

\begin{center}
{\Large Supplementary material}\\[4pt]
{\large How Much of a Time Series Survives Its Image? A Cross-Domain Comparative Analysis of the Reconstruction Capacity of Image-Based Representations}
\end{center}

\noindent Section, table, figure and equation numbers without the prefix S refer to the main text. Part A describes the experimental set-up, Part B proves Proposition~\ref{M-prop:rp} and checks its hypotheses on the data, and Part C extends the results of the main text in the order of its research questions. Every table except Table~\ref{tab:s_storage} is generated from the stored results by the scripts named in the text, which are available with the code.

% ============================================================================================
\partheading{Part A. Experimental set-up}

\section{Datasets and preprocessing}\label{sec:s_datasets}
All datasets are public. Each recording is cut into non-overlapping windows of $T$ samples that never cross a label change, a time gap or a missing value; windows that would are dropped. Every channel is a separate series, and series whose range is below $10^{-12}$ are excluded as constant. Table~\ref{tab:s_samples} gives the resulting sample sizes.

\begin{itemize}
\item \textbf{WISDM}~\cite{misc_wisdm_smartphone_and_smartwatch_activity_and_biometrics_dataset__507}. Smartwatch accelerometer ($x$, $y$, $z$, m/s$^2$) at 20\,Hz, 51 subjects and five activities (walking, jogging, stairs, sitting, standing). One recording per maximal run of constant subject and activity; timestamps are not used. Group: subject.
\item \textbf{PAMAP2}~\cite{reiss2012pamap2}. The $\pm16$\,g accelerometers of the three inertial units (hand, chest, ankle; nine channels) at 100\,Hz, nine subjects and the twelve protocol activities. Transient periods (activity 0), the heart rate and the saturating $\pm6$\,g accelerometers are not used. Gaps of up to ten samples are filled by linear interpolation; windows with longer gaps are dropped. Group: subject.
\item \textbf{ESC-50}~\cite{piczak2015esc}. 2\,000 five-second clips in 50 classes, mono 16-bit PCM scaled to $[-1,1]$ and resampled from 44.1 to 8\,kHz with an anti-aliased polyphase filter. Group: the Freesound source recording, so clips cut from one recording never fall on both sides of a split. The audio study uses the official ESC-10 subset (400 clips) at 2, 4 and 8\,kHz and removes windows with a standard deviation below $10^{-3}$, which are silent padding.
\item \textbf{PTB-XL}~\cite{wagner2020ptbxl}. Ten-second 12-lead electrocardiograms in mV, read from the 100\,Hz files with each lead's gain and baseline; 3\,000 of the 21\,799 records are drawn at random (seed 0). Labels are the five diagnostic superclasses, as a multilabel task. Group: patient. The sampling-rate study reads the same records from the 500\,Hz files.
\item \textbf{Swarm}~\cite{friis2006swarm}. Vector magnetic field $B_{\mathrm{NEC}}$ (nT) of the MAGx\_LR\_1B product at 1\,Hz, satellites A, B and C, 1--31 January 2023. Samples whose quality flag is 255 or has bit 2 or 16 set are treated as missing. Unlabelled; group: day.
\item \textbf{STEAD}~\cite{mousavi2019stead}. Three-component seismograms (E, N, Z) of 60\,s at 100\,Hz from the first two data chunks; a class-balanced random subset of 1\,000 traces per class (local earthquake, noise) is drawn with seed 0. Group: the earthquake event, or the station for noise, so neither an event nor a noise station falls on both sides of a split.
\item \textbf{EMG}~\cite{krilova2018emg}. Eight-channel armband surface EMG, 36 subjects, static gestures of 3\,s. Pauses (label 0) and the extended-palm gesture (label 7, performed by few subjects) are excluded, leaving six gestures. The logged time advances in 1\,ms steps by repeating values while the signal changes about 100 times per second, so each file is resampled by sample-and-hold onto a uniform 100\,Hz grid. Group: subject.
\end{itemize}

\begin{table}[H]
\centering
\caption{Sample sizes. Fidelity: windows drawn per dataset and length (seed 0) and the non-constant channel series they contain. Classification: windows drawn per dataset and length; ESC-50 is evaluated as ESC-10 instead (Table~\ref{tab:s_esc10}) and Swarm is unlabelled. ``--'': fewer than 200 windows available.}\label{tab:s_samples}
\footnotesize
\setlength{\tabcolsep}{4pt}
\begin{tabular}{@{}lrrrrrr@{}}
\toprule
& \multicolumn{3}{c}{Fidelity: windows / channel series} & \multicolumn{3}{c}{Classification: windows} \\ \cmidrule(lr){2-4}\cmidrule(l){5-7}
Dataset & $T=64$ & $T=128$ & $T=256$ & $T=64$ & $T=128$ & $T=256$ \\
\midrule
WISDM & 500 / 1500 & 500 / 1500 & 500 / 1500 & 6000 & 6000 & 4073 \\
PAMAP2 & 500 / 4500 & 500 / 4500 & 500 / 4500 & 6000 & 6000 & 6000 \\
ESC-50 & 3000 / 2704 & 3000 / 2704 & 3000 / 2699 & -- & -- & -- \\
PTB-XL & 500 / 6000 & 500 / 6000 & 500 / 6000 & 6000 & 6000 & 6000 \\
Swarm & 500 / 1500 & 500 / 1500 & 500 / 1500 & -- & -- & -- \\
STEAD & 500 / 1494 & 500 / 1488 & 500 / 1490 & 6000 & 6000 & 6000 \\
EMG & 500 / 4000 & 500 / 4000 & -- & 1964 & 880 & -- \\
\bottomrule
\end{tabular}

\end{table}

\section{Storage and image-only decoding}\label{sec:s_storage}
Table~\ref{tab:s_storage} lists, for every representation, the normalised domain of the series, the 8-bit storage rule $\mathcal Q_m$ and what the decoder recovers from the image before the calibration is applied (Section~\ref{M-sec:storage}). Fixed storage ranges depend only on the method; the per-image ranges of the RP are never passed to a decoder.

\begin{table}[H]
\centering
\caption{Normalisation and 8-bit storage of each representation. Rounding is to the nearest of 256 levels.}\label{tab:s_storage}
\footnotesize
\setlength{\tabcolsep}{4pt}
\begin{tabular}{@{}L{0.17\textwidth}L{0.2\textwidth}L{0.29\textwidth}L{0.26\textwidth}@{}}
\toprule
Representation & Series domain & Stored image & Decoder input \\
\midrule
Signed RP (conventional and circular) & window as recorded & min--max of the matrix, per image (contains the zero diagonal) & matrix re-centred at its diagonal level, up to a positive factor \\
Unsigned RP (conventional and circular) & window as recorded & divided by its maximum & distances up to a positive factor \\
GASF & $[0,1]$ & fixed range $[-1,1]$ & matrix \\
GADF & $[-0.7,0.7]$ & fixed range $[-1,1]$ & matrix \\
MTF & quantile states, $Q=4$ & fixed range $[0,1]$ & matrix \\
XIRP & $[1,2]$ & fixed ranges: $[1,2]$ on the diagonal, $[-\log2,\log2]$ off it & matrix \\
\bottomrule
\end{tabular}
\end{table}

Every decoder receives the stored image and the channel calibration $(x_{\max},x_{\min})$, never the series. Whether the calibration is used for anything but the final map to physical units can be tested: if it chose a reflection, a branch or a threshold, decoding the same image with another calibration would not give an affine transform of the first decoding. The script \texttt{\_check\_image\_only.py} decodes the 8-bit images of \imgOnlySeries{} series with their own calibration and with an unrelated one, for all ten combinations; the largest deviation from the affine transform is \imgOnlyMax{} of the range. One of the series has no sign-informative pair, so its two reflections tie exactly. In an earlier version, the three signed decoders broke that tie after mapping the candidates to physical units, so the rounding of the calibration decided it; they now compare the candidates on the unit range. On the benchmark windows, where no exact tie occurs, the change moves the reconstructions by less than $10^{-13}$ of the range, below the precision of every reported number.

\section{Fixed method parameters}\label{sec:s_params}
No parameter of the encoders or decoders is tuned on the evaluated windows except where stated. The GADF interval $[-0.7,0.7]$ follows from the branch condition $\alpha<1/\sqrt2$ (Section~\ref{M-sec:formulation}). The number of GADF pivots (five) and of MTF states ($Q=4$) were fixed in our earlier single-dataset study on WISDM, which is also one of the seven datasets here, and were not changed for this benchmark. The MTF encoder uses $Q=4$ quantile states; the decoder uses the same $Q$. It normalises the image rows to unit length and clusters them by $k$-means with cosine similarity, initialised by farthest-point sampling from a fixed seed (0), for at most 50 Lloyd iterations; an empty cluster is re-seeded with the row farthest from every centroid. The transition matrix $\hat W$ is the mean of each block of the image, normalised by rows, and the clusters are sorted by the Fiedler vector of the unnormalised Laplacian of $(\hat W+\hat W^\top)/2$ with the sign returned by the eigensolver. Hirata's recurrence rate ($0.75$) was fixed before the benchmark; its sensitivity is reported in Section~\ref{sec:s_hirata}. Sipers' frequency threshold ($10^{-4}$) was first chosen in an exploratory sweep on evaluated windows (Section~\ref{sec:s_sipers}) and confirmed on calibration windows outside the evaluated samples (Section~\ref{sec:s_calibration}).

\section{Classifier settings}\label{sec:s_classifiers}
All classifiers are trained on original windows, standardised per channel with the statistics of the training fold, and evaluated on the original and on every reconstruction of the same held-out windows. KNN ($k=5$), an RBF SVM ($C=1$, Platt probabilities, one-vs-rest for the multilabel PTB-XL) and a random forest (300 trees) take the flattened window. The LSTM and GRU (64 and 32 units, dropout $0.2$) and a one-dimensional CNN (convolutions of 32 and 64 filters of width 5, max pooling, global average pooling and a 64-unit dense layer) take the $(T,C)$ window and are trained with Adam, batch size 64 and at most 20 epochs, with early stopping (patience 4) on the last 10\% of the training fold. No hyperparameter is tuned. Metrics are computed per fold and averaged over the three folds; the equivalence intervals resample the groups of the split.

% ============================================================================================
\partheading{Part B. Proposition~\ref{M-prop:rp}}

\section{Proof of Proposition~\ref{M-prop:rp}}\label{sec:s_proof}
We use the notation of Section~\ref{M-sec:rp_encodings} and the decoder steps of Section~\ref{M-sec:rp_decoder} of the main text. Let $S\in\mathbb R^{(T-1)\times2}$ have rows $s_i=(x_i,x_{i+1})$, and write $O(2)$ for the orthogonal $2\times2$ matrices. The proof follows the decoder: geometry (step~2), overlap (step~3), translation and scale (step~4), and sign (step~5).

\paragraph{Step 1: geometry}
By (a), the centred states have affine rank two, so classical multidimensional scaling applied to the exact distance matrix $|R|$ returns
\begin{equation}
Y=SP_0+\mathbf 1t^\top,\qquad P_0\in O(2),\ t\in\mathbb R^2~\cite{Torgerson1952}.
\end{equation}
Every candidate configuration produced by step~3 applies a further orthogonal transform to $Y$, so it is of the form $Z=SP+\mathbf 1c^\top$ with $P\in O(2)$ and $c\in\mathbb R^2$.

\paragraph{Step 2: overlap}
Write $(a_i,b_i)$ for the rows of $Z$, so that $a_i=P_{11}x_i+P_{21}x_{i+1}+c_1$ and $b_i=P_{12}x_i+P_{22}x_{i+1}+c_2$. Then
\begin{equation}\label{eq:s_residual}
b_i-a_{i+1}=P_{12}\,x_i+(P_{22}-P_{11})\,x_{i+1}-P_{21}\,x_{i+2}+(c_2-c_1),
\end{equation}
so the residual is constant in $i$ if and only if $P_{12}x_i+(P_{22}-P_{11})x_{i+1}-P_{21}x_{i+2}$ is. Every $P\in O(2)$ is a rotation or a reflection, with $c^2+s^2=1$:
\begin{align}
P&=\begin{pmatrix}c&s\\-s&c\end{pmatrix}:&& b_i-a_{i+1}=s\,(x_i+x_{i+2})+\text{const},\label{eq:s_rot}\\
P&=\begin{pmatrix}c&s\\s&-c\end{pmatrix}:&& b_i-a_{i+1}=s\,(x_i-x_{i+2})-2c\,x_{i+1}+\text{const}.\label{eq:s_ref}
\end{align}
Under (b), relation (R) fails, so~\eqref{eq:s_rot} is constant only for $s=0$, that is $P=\pm I$. Relation (F) also fails, and~\eqref{eq:s_ref} constant would be (F) with $\alpha=s$ and $\beta=-2c$, where $(\alpha,\beta)\neq(0,0)$ because $(s,c)$ lies on the unit circle; so no reflection satisfies the overlap. Hence the only configurations that satisfy the overlap exactly are $P=\pm I$, i.e.\ the states of $x$ and of $-x$, up to a translation.

\paragraph{Step 3: the matrices $H_\eta$}
In step~3 of the decoder, the transforms of parity $\eta$ that satisfy the overlap exactly are the unit vectors $u$ with $u^\top H_\eta u=0$. By Step~2 there are exactly two, $\pm u_0$, of one parity, and none of the other. Therefore one $H_\eta$ has a zero eigenvalue whose eigenvector, with both signs, yields the candidates $x$ and $-x$. That eigenvalue is simple, because a vanishing $H_\eta$ would make every transform of that parity a solution, and the other matrix, which is positive semidefinite with no zero eigenvalue, is positive definite.

Conversely, if (R) holds, every rotation satisfies~\eqref{eq:s_rot} and the corresponding $H_\eta$ vanishes; if (F) holds, a reflection satisfies~\eqref{eq:s_ref}, giving a zero eigenvalue of the other parity. Under (a), condition (b) is therefore equivalent to the eigenvalue pattern stated in the proposition, which can be checked from the image.

If the centred triples $(x_i,x_{i+1},x_{i+2})$ have rank three, no nontrivial affine relation among three consecutive samples holds for all $i$, so neither (R) nor (F) can hold; this is the sufficient criterion of the proposition.

\paragraph{Step 4: translation and scale}
Step~4 of the decoder removes the residual translation $\delta^\star$ and assembles the samples of every candidate; the final range map removes the remaining translation and positive scale. The two candidates of the parity that satisfies the overlap are therefore $\pm x$ on the normalised domain.

\paragraph{Step 5: sign}
In step~5 the assembled series of every candidate is re-encoded, and the candidate $x$ reproduces the observed matrix exactly after both are normalised to unit Frobenius norm. Any candidate whose signed RP also reproduces it has states whose distance matrix is proportional to $|R|$ and, being the states of a series, satisfies the overlap by construction, so it is $\pm x$ up to scale and translation by Steps~1--4. Reflecting the series maps $s_i-s_j$ to $s_j-s_i$, so $\sigma_{ij}$ becomes $\sigma_{ji}$ in~\eqref{M-eq:signed_rp}: on a sign-informative pair the entry changes sign, and on a non-informative pair it does not. The signed RPs of $x$ and $-x$ therefore differ if and only if (c) holds. Under (c) the selected candidate is $x$; without (c) the two are indistinguishable from the image. $\square$

\section{Hypotheses on the benchmark series}\label{sec:s_hypotheses}
The script \texttt{\_check\_sign\_identifiability.py} evaluates the hypotheses of Proposition~\ref{M-prop:rp} on every channel series of the benchmark, with a relative tolerance of $10^{-9}$ and the same exclusion of constant series as the benchmark. Each hypothesis is measured by a scale-free margin, and holds when its margin exceeds the tolerance:
\begin{itemize}
\item (a): the ratio $\sigma_2/\sigma_1$ of the singular values of the centred states;
\item (R): the range of $x_i+x_{i+2}$ divided by the range of $x$;
\item (F): the ratio of the smallest to the largest singular value of the column-normalised matrix with columns $x_{i+2}-x_i$, $x_{i+1}$ and $1$;
\item rank three: the ratio of the smallest to the largest singular value of the centred triples;
\item (c): the share of state pairs whose displacement lies within $\pi/4$ of $\pm(1,1)$.
\end{itemize}
Table~\ref{tab:s_hypotheses} gives the smallest margin per dataset over its three lengths. No series fails any hypothesis. The smallest margins, on Swarm for (a) and on ESC-50 for (F), are those of nearly collinear states, which is the regime where storage error grows (Section~\ref{M-sec:results_conditions}). The check is numerical evidence that the benchmark meets the hypotheses, not a substitute for the proof. The same script confirms, with the production encoder, that the alternating series of the main text and its negative have identical signed RPs.

\begin{table}[H]
\centering
\caption{Hypotheses of Proposition~\ref{M-prop:rp} on every benchmark series: smallest margin per dataset over $T\in\{64,128,256\}$, and number of series that fail any of (a)--(c).}\label{tab:s_hypotheses}
\footnotesize
\setlength{\tabcolsep}{4pt}
\resizebox{\textwidth}{!}{\begin{tabular}{@{}lrrrrrrr@{}}
\toprule
Dataset & Series & (a) & (R) & (F) & Rank 3 & Sign-inf. share & Failures \\
\midrule
WISDM & 4\,500 & 2.3e-02 & 5.5e-01 & 7.2e-04 & 9.9e-03 & 0.230 & 0 \\
PAMAP2 & 13\,500 & 2.3e-02 & 6.8e-01 & 2.4e-03 & 9.4e-03 & 0.360 & 0 \\
ESC-50 & 8\,107 & 6.7e-03 & 1.0e-01 & 3.1e-06 & 1.1e-03 & 0.064 & 0 \\
PTB-XL & 18\,000 & 3.4e-02 & 6.7e-01 & 1.9e-02 & 8.5e-03 & 0.476 & 0 \\
Swarm & 4\,500 & 5.7e-06 & 1.3e+00 & 2.2e-05 & 1.2e-06 & 0.830 & 0 \\
STEAD & 4\,472 & 7.0e-03 & 5.0e-01 & 9.1e-02 & 1.9e-05 & 0.213 & 0 \\
EMG & 8\,000 & 3.1e-01 & 7.7e-01 & 3.3e-01 & 2.2e-01 & 0.231 & 0 \\
\bottomrule
\end{tabular}
}
\end{table}

% ============================================================================================
\partheading{Part C. Additional results}

\section{Fidelity at the other window lengths (RQ1)}\label{sec:s_lengths}
Tables~\ref{tab:s_l64} and~\ref{tab:s_l256} repeat Table~\ref{M-tab:fidelity} of the main text for $T=64$ and $T=256$. EMG at $T=256$ has fewer than 200 windows and is not reported.

\begin{table}[H]
\centering
\caption{Mean 8-bit NMAE of every encoding--decoder combination at $T=64$. Uns./Sig.: unsigned and signed RP.}\label{tab:s_l64}
\scriptsize
\setlength{\tabcolsep}{2.5pt}
\begin{tabular}{@{}lrrrrrrrrrrr@{}}
\toprule
& \multicolumn{2}{c}{Hirata} & \multicolumn{2}{c}{Sipers} & \multicolumn{3}{c}{Geometric (ours)} & & & & \\
\cmidrule(lr){2-3}\cmidrule(lr){4-5}\cmidrule(lr){6-8}
Dataset & Uns. & Sig. & Uns. & Sig. & Uns. & Oracle & Sig. & GASF & GADF & MTF & XIRP \\
\midrule
WISDM & 0.2647 & 0.1171 & 0.2084 & 0.0042 & 0.2039 & 0.0005 & 0.0012 & 0.0016 & 0.0011 & 0.2930 & 0.0003 \\
PAMAP2 & 0.2630 & 0.1153 & 0.2108 & 0.0049 & 0.2103 & 0.0008 & 0.0017 & 0.0017 & 0.0012 & 0.2889 & 0.0003 \\
ESC-50 & 0.2463 & 0.0956 & 0.1997 & 0.0056 & 0.2247 & 0.0007 & 0.0016 & 0.0017 & 0.0010 & 0.2937 & 0.0003 \\
PTB-XL & 0.3530 & 0.1916 & 0.2991 & 0.0060 & 0.2910 & 0.0010 & 0.0022 & 0.0033 & 0.0019 & 0.3434 & 0.0003 \\
Swarm & 0.2710 & 0.0293 & 0.2613 & 0.0028 & 0.2630 & 0.0129 & 0.0144 & 0.0025 & 0.0013 & 0.2942 & 0.0003 \\
STEAD & 0.2361 & 0.1001 & 0.2052 & 0.0055 & 0.2030 & 0.0006 & 0.0014 & 0.0016 & 0.0010 & 0.2733 & 0.0003 \\
EMG & 0.2300 & 0.1365 & 0.1869 & 0.0039 & 0.1947 & 0.0004 & 0.0010 & 0.0014 & 0.0020 & 0.2982 & 0.0006 \\
\bottomrule
\end{tabular}

\end{table}

\begin{table}[H]
\centering
\caption{As Table~\ref{tab:s_l64}, at $T=256$.}\label{tab:s_l256}
\scriptsize
\setlength{\tabcolsep}{2.5pt}
\begin{tabular}{@{}lrrrrrrrrrr@{}}
\toprule
& \multicolumn{2}{c}{Hirata} & \multicolumn{2}{c}{Sipers} & \multicolumn{2}{c}{Geometric (ours)} & & & & \\
\cmidrule(lr){2-3}\cmidrule(lr){4-5}\cmidrule(lr){6-7}
Dataset & Uns. & Sig. & Uns. & Sig. & Uns. & Sig. & GASF & GADF & MTF & XIRP \\
\midrule
WISDM & 0.2228 & 0.0971 & 0.1743 & 0.0046 & 0.1682 & 0.0010 & 0.0014 & 0.0015 & 0.2700 & 0.0002 \\
PAMAP2 & 0.2084 & 0.0960 & 0.1829 & 0.0053 & 0.1736 & 0.0015 & 0.0015 & 0.0015 & 0.2740 & 0.0002 \\
ESC-50 & 0.1987 & 0.0755 & 0.1747 & 0.0058 & 0.1759 & 0.0017 & 0.0015 & 0.0011 & 0.2672 & 0.0001 \\
PTB-XL & 0.3066 & 0.1600 & 0.2588 & 0.0051 & 0.2716 & 0.0012 & 0.0025 & 0.0026 & 0.3291 & 0.0002 \\
Swarm & 0.2950 & 0.0339 & 0.3111 & 0.0155 & 0.2573 & 0.0174 & 0.0027 & 0.0009 & 0.3066 & 0.0001 \\
STEAD & 0.1723 & 0.0789 & 0.1581 & 0.0047 & 0.1552 & 0.0012 & 0.0014 & 0.0011 & 0.2592 & 0.0001 \\
EMG & -- & -- & -- & -- & -- & -- & -- & -- & -- & -- \\
\bottomrule
\end{tabular}

\end{table}

\section{Spread of the per-series error (RQ1)}\label{sec:s_spread}
Table~\ref{M-tab:fidelity} of the main text reports means, which a few poorly conditioned series can dominate. Table~\ref{tab:s_spread} gives the median and the interquartile range of the same per-series errors. For the unsigned RP the distribution is bimodal, according to whether the reflection is returned, so its median is not representative.

\begin{table}[H]
\centering
\caption{Median [interquartile range] of the per-series 8-bit NMAE at $T=128$, on the series of Table~\ref{M-tab:fidelity}.}\label{tab:s_spread}
\scriptsize
\setlength{\tabcolsep}{2.5pt}
\resizebox{\textwidth}{!}{\begin{tabular}{@{}lrrrrrrr@{}}
\toprule
Decoder & WISDM & PAMAP2 & ESC-50 & PTB-XL & Swarm & STEAD & EMG \\
\midrule
Hirata\,/\,uns. & 0.1850 [0.2753] & 0.2093 [0.2632] & 0.1625 [0.2648] & 0.2848 [0.3557] & 0.0979 [0.4861] & 0.1478 [0.2211] & 0.1857 [0.1620] \\
Hirata\,/\,sig. & 0.0901 [0.0463] & 0.0829 [0.0549] & 0.0764 [0.0381] & 0.1420 [0.2065] & 0.0180 [0.0220] & 0.0837 [0.0335] & 0.1101 [0.0341] \\
Sipers\,/\,uns. & 0.1702 [0.3670] & 0.0114 [0.3569] & 0.1782 [0.3603] & 0.0076 [0.5704] & 0.5025 [0.5056] & 0.2399 [0.3343] & 0.0099 [0.3319] \\
Sipers\,/\,sig. & 0.0037 [0.0033] & 0.0045 [0.0046] & 0.0055 [0.0040] & 0.0047 [0.0030] & 0.0030 [0.0029] & 0.0044 [0.0032] & 0.0031 [0.0025] \\
Ours\,/\,uns. & 0.1614 [0.3726] & 0.1138 [0.3598] & 0.2372 [0.3624] & 0.1669 [0.5936] & 0.0184 [0.4911] & 0.2141 [0.3385] & 0.1964 [0.3316] \\
Ours\,/\,sig. & 0.0008 [0.0004] & 0.0011 [0.0008] & 0.0009 [0.0007] & 0.0014 [0.0007] & 0.0133 [0.0074] & 0.0008 [0.0004] & 0.0008 [0.0002] \\
GASF & 0.0013 [0.0007] & 0.0014 [0.0008] & 0.0015 [0.0006] & 0.0015 [0.0028] & 0.0024 [0.0005] & 0.0013 [0.0005] & 0.0012 [0.0006] \\
GADF & 0.0010 [0.0008] & 0.0011 [0.0009] & 0.0009 [0.0006] & 0.0016 [0.0021] & 0.0013 [0.0014] & 0.0009 [0.0006] & 0.0020 [0.0014] \\
MTF & 0.2903 [0.2599] & 0.2851 [0.2665] & 0.2863 [0.3132] & 0.3505 [0.1182] & 0.4330 [0.4366] & 0.2673 [0.2972] & 0.2769 [0.0967] \\
XIRP & 0.0002 [0.0001] & 0.0002 [0.0001] & 0.0002 [0.0001] & 0.0002 [0.0001] & 0.0002 [0.0001] & 0.0002 [0.0001] & 0.0004 [0.0004] \\
\bottomrule
\end{tabular}
}
\end{table}

\section{Spectral and correlational fidelity (RQ1)}\label{sec:s_spectral}
Table~\ref{tab:s_spectral} separates the two reasons for a large pointwise error discussed in Section~\ref{M-sec:results_fidelity}: decoders that return the reflection keep the spectrum and the energy, while those that lose the waveform distort them.

\begin{table}[H]
\centering
\caption{Spectral and correlational fidelity after 8-bit storage: range of the dataset means over the seven datasets at $T=128$. Pearson: correlation with the original; PSD shape: $L_1$ distance between unit-mass Welch spectra; Energy: relative error of total power.}\label{tab:s_spectral}
\footnotesize
\begin{tabular}{@{}lrrr@{}}
\toprule
Decoder & Pearson & PSD shape & Energy \\
\midrule
Hirata\,/\,uns. & -0.023--0.019 & 0.0550--0.6117 & 0.224--0.766 \\
Hirata\,/\,sig. & 0.556--0.982 & 0.0562--0.6053 & 0.225--0.730 \\
Sipers\,/\,uns. & -0.075--0.038 & 0.0021--0.0125 & 0.005--0.012 \\
Sipers\,/\,sig. & 0.999--1.000 & 0.0063--0.0260 & 0.010--0.020 \\
Ours\,/\,uns. & -0.036--0.045 & 0.0007--0.0151 & 0.001--0.036 \\
Ours\,/\,sig. & 0.999--1.000 & 0.0034--0.0187 & 0.002--0.055 \\
GASF & 0.999--1.000 & 0.0041--0.0171 & 0.005--0.012 \\
GADF & 1.000--1.000 & 0.0005--0.0048 & 0.001--0.002 \\
MTF & -0.021--0.009 & 0.3279--0.7871 & 1.053--3.154 \\
XIRP & 1.000--1.000 & 0.0003--0.0017 & 0.000--0.001 \\
\bottomrule
\end{tabular}

\end{table}

\section{Recurrence-rate sensitivity of Hirata's reconstruction (RQ2)}\label{sec:s_hirata}
Table~\ref{tab:s_hirata} gives the 8-bit NMAE of Hirata's reconstruction~\cite{hirata2008reproduction} against the recurrence rate used to threshold the image ($T=128$, 100 benchmark windows per dataset, every channel). The rate of $0.75$ used in the main text was fixed in advance; choosing the best rate per dataset with knowledge of the result lowers the error by at most \hirataGainUns{} on the unsigned and \hirataGainSig{} on the signed RP. Because this sweep uses evaluated windows, it is reported as a sensitivity analysis, not used to choose the rate.

\begin{table}[H]
\centering
\caption{Hirata's reconstruction: 8-bit NMAE against the recurrence rate.}\label{tab:s_hirata}
\scriptsize
\setlength{\tabcolsep}{3pt}
\begin{tabular}{@{}llrrrrrrrrr@{}}
\toprule
Dataset & Encoding & 0.02 & 0.05 & 0.1 & 0.2 & 0.3 & 0.5 & 0.7 & 0.75 & Best \\
\midrule
WISDM & unsigned & 0.239 & 0.250 & 0.248 & 0.245 & 0.239 & 0.240 & 0.237 & 0.231 & 0.75 \\
 & signed & 0.225 & 0.196 & 0.171 & 0.136 & 0.120 & 0.103 & 0.104 & 0.103 & 0.75 \\
PAMAP2 & unsigned & 0.240 & 0.255 & 0.252 & 0.238 & 0.236 & 0.231 & 0.233 & 0.238 & 0.5 \\
 & signed & 0.222 & 0.194 & 0.161 & 0.124 & 0.105 & 0.098 & 0.102 & 0.107 & 0.5 \\
ESC-50 & unsigned & 0.237 & 0.253 & 0.242 & 0.251 & 0.236 & 0.197 & 0.215 & 0.215 & 0.5 \\
 & signed & 0.217 & 0.156 & 0.115 & 0.088 & 0.078 & 0.076 & 0.079 & 0.081 & 0.5 \\
PTB-XL & unsigned & 0.307 & 0.311 & 0.320 & 0.323 & 0.326 & 0.330 & 0.342 & 0.328 & 0.02 \\
 & signed & 0.290 & 0.287 & 0.285 & 0.277 & 0.272 & 0.252 & 0.213 & 0.193 & 0.75 \\
Swarm & unsigned & 0.273 & 0.265 & 0.287 & 0.291 & 0.293 & 0.296 & 0.293 & 0.290 & 0.05 \\
 & signed & 0.064 & 0.033 & 0.021 & 0.020 & 0.024 & 0.028 & 0.028 & 0.028 & 0.2 \\
STEAD & unsigned & 0.239 & 0.251 & 0.235 & 0.222 & 0.224 & 0.219 & 0.221 & 0.212 & 0.75 \\
 & signed & 0.216 & 0.162 & 0.124 & 0.096 & 0.085 & 0.079 & 0.084 & 0.085 & 0.5 \\
EMG & unsigned & 0.214 & 0.234 & 0.234 & 0.220 & 0.215 & 0.201 & 0.200 & 0.203 & 0.7 \\
 & signed & 0.199 & 0.192 & 0.163 & 0.138 & 0.126 & 0.116 & 0.119 & 0.121 & 0.5 \\
\bottomrule
\end{tabular}

\end{table}

\section{Sipers et al.\ and the geometric decoder on the same image (RQ2)}\label{sec:s_sipers}
The reconstruction of Sipers et al.~\cite{sipers2017robust} assumes a periodic signal and a circular delay embedding, so in the main benchmark it decodes the circular variant of the RP, while the geometric decoder decodes the conventional RP. The circular RP is built on the states $(x_i,x_{i+1})$, $i=1,\ldots,T$, with $x_{T+1}=x_1$: it is the conventional RP with one added state, the wrap state $(x_T,x_1)$, stored with the same 8-bit maps. Sipers' method represents the window by its periodic extension, which a finite sequence always admits, so its hypotheses hold for the circular embedding of any window; the cost is the wrap state, which joins the last sample to the first and is not part of the recorded dynamics.

Table~\ref{tab:s_sipers} removes the difference of encoding: both decoders receive the identical circular image of the first 100 benchmark windows of each dataset at $T=128$, every channel, with the same per-channel calibration, and on the signed RP both choose the reflection from the image alone with the same criterion, so any gap is due to the decoder. The unsigned columns fix the reflection with the original series; they are a bound on what the geometry alone allows, not a reconstruction.

Sipers' decoder discards frequencies whose power, estimated from the image, is below a share of the largest; this frequency threshold does not use the original series. With a threshold of zero the implementation is exact in floating point (median error $0$; maximum \sipZeroMaxClean{} on WISDM, Swarm and STEAD, and \sipZeroMaxPtb{} for the worst of \sipZeroPtbSeries{} PTB-XL series, where near-zero Fourier coefficients enter the division of Eq.~(15) of Sipers et al.). With the threshold of $0.01$ used by its authors, the error is \sipPaperTolMin{} to \sipPaperTolMax{} even in floating point, because our windows are broadband. The last columns of Table~\ref{tab:s_sipers} sweep the threshold on these evaluated windows. This exploratory sweep is where the value used throughout, $10^{-4}$, was first chosen, so it is not independent of the evaluation; Section~\ref{sec:s_calibration} confirms it on separate calibration windows.

\begin{table}[H]
\centering
\caption{Mean 8-bit NMAE on the same circular RP ($T=128$, first 100 benchmark windows per dataset). Unsigned: reflection fixed with the original series, a bound on the geometry alone. Signed: reflection chosen from the image by the sign, as in the main text. Last columns: Sipers' decoder on the unsigned image against the frequency threshold.}\label{tab:s_sipers}
\scriptsize
\setlength{\tabcolsep}{3pt}
\begin{tabular}{@{}lrrrrrrrrr@{}}
\toprule
& \multicolumn{2}{c}{Unsigned, oracle sign} & \multicolumn{2}{c}{Signed} & \multicolumn{5}{c}{Sipers, unsigned, oracle sign, node threshold} \\
\cmidrule(lr){2-3}\cmidrule(lr){4-5}\cmidrule(l){6-10}
Dataset & Sipers & Geometric & Sipers & Geometric & 0.01 & 0.001 & 0.0001 & 1e-06 & 0 \\
\midrule
WISDM & 0.0026 & 0.0004 & 0.0042 & 0.0010 & 0.0440 & 0.0092 & 0.0026 & 0.0031 & 0.0034 \\
PAMAP2 & 0.0033 & 0.0006 & 0.0050 & 0.0014 & 0.0296 & 0.0082 & 0.0033 & 0.0042 & 0.0043 \\
ESC-50 & 0.0041 & 0.0006 & 0.0054 & 0.0013 & 0.0343 & 0.0099 & 0.0041 & 0.0049 & 0.0049 \\
PTB-XL & 0.0035 & 0.0007 & 0.0054 & 0.0014 & 0.0216 & 0.0052 & 0.0035 & 0.0050 & 0.0054 \\
Swarm & 0.0019 & 0.0010 & 0.0046 & 0.0025 & 0.0410 & 0.0379 & 0.0019 & 0.0020 & 0.0020 \\
STEAD & 0.0040 & 0.0005 & 0.0054 & 0.0012 & 0.0271 & 0.0078 & 0.0040 & 0.0055 & 0.0061 \\
EMG & 0.0015 & 0.0003 & 0.0039 & 0.0008 & 0.0093 & 0.0016 & 0.0015 & 0.0017 & 0.0017 \\
\bottomrule
\end{tabular}

\end{table}

\section{Calibration of Sipers' frequency threshold (RQ2)}\label{sec:s_calibration}
The threshold first chosen in Section~\ref{sec:s_sipers} is checked on \calWindows{} windows per dataset at $T=128$ that lie outside both the fidelity and the classification samples, drawn with a different seed (seed 1; script \texttt{\_sipers\_calibration.py}). Such windows exist for \calDatasets{} of the seven datasets; all EMG windows are evaluated. Averaging the 8-bit NMAE over datasets with equal weight, the threshold with the lowest error is \calTolOracle{} when the reflection is fixed with the original series and \calTolSigned{} when it is chosen from the signed image, the value used throughout. The calibration windows may come from the same recordings as evaluated windows.

\section{Sampling rate at a fixed length and at a fixed duration (RQ3)}\label{sec:s_sampling}
Table~\ref{tab:s_sampling} separates the sampling rate from the window duration on PTB-XL. With 128 samples, a higher rate also means a shorter window; with a fixed duration of 1.28\,s, the number of samples grows with the rate. In both settings the conditioning $\sigma_2/\sigma_1$ falls and the error of the signed RP rises with the rate, while GASF, GADF and XIRP do not follow that trend. The ``native'' rows use the 100\,Hz files distributed with the dataset instead of resampling the 500\,Hz records.

\begin{table}[H]
\centering
\caption{PTB-XL, anti-aliased resampling from 500\,Hz: median $\sigma_2/\sigma_1$ of the RP states, share of series below $0.1$, and mean 8-bit NMAE.}\label{tab:s_sampling}
\scriptsize
\setlength{\tabcolsep}{1.8pt}
\begin{tabular}{@{}llrrrrrrrr@{}}
\toprule
Window & Rate & Seconds & $\sigma_2/\sigma_1$ & Share $<0.1$ & Signed RP & Sipers\,/\,sig. & GASF & GADF & XIRP \\
\midrule
128 samples & 500 Hz & 0.26 & 0.08 & 67\% & 0.0076 & 0.0078 & 0.0032 & 0.0018 & 0.0005 \\
 & 250 Hz & 0.51 & 0.16 & 23\% & 0.0041 & 0.0074 & 0.0032 & 0.0019 & 0.0002 \\
 & 100 Hz & 1.28 & 0.37 & 0\% & 0.0015 & 0.0055 & 0.0029 & 0.0022 & 0.0002 \\
 & 100 (native) Hz & 1.28 & 0.37 & 0\% & 0.0015 & 0.0054 & 0.0028 & 0.0022 & 0.0002 \\
 & 50 Hz & 2.56 & 0.61 & 0\% & 0.0009 & 0.0040 & 0.0021 & 0.0019 & 0.0002 \\
\midrule
1.28 s & 500 Hz & 1.28 & 0.08 & 76\% & 0.0079 & 0.0061 & 0.0028 & 0.0032 & 0.0002 \\
 & 250 Hz & 1.28 & 0.16 & 9\% & 0.0026 & 0.0064 & 0.0028 & 0.0028 & 0.0002 \\
 & 100 Hz & 1.28 & 0.37 & 0\% & 0.0015 & 0.0055 & 0.0029 & 0.0022 & 0.0002 \\
 & 100 (native) Hz & 1.28 & 0.37 & 0\% & 0.0015 & 0.0054 & 0.0028 & 0.0022 & 0.0002 \\
 & 50 Hz & 1.28 & 0.64 & 0\% & 0.0012 & 0.0037 & 0.0024 & 0.0016 & 0.0003 \\
\bottomrule
\end{tabular}

\end{table}

\section{Robustness per dataset (RQ3)}\label{sec:s_robustness}
Table~\ref{tab:s_robustness} gives, for each dataset, the clean 8-bit error and the error with Gaussian noise of $\sigma=16$ intensity levels on the stored image, which Figure~\ref{M-fig:robustness_sampling}a averages over datasets.

\begin{table}[H]
\centering
\caption{Mean 8-bit NMAE at $T=128$: clean / with noise of $\sigma=16$ (200 windows per dataset, two replicates). Sipers' decoder reads the circular RP.}\label{tab:s_robustness}
\scriptsize
\setlength{\tabcolsep}{3pt}
\begin{tabular}{@{}lrrrrr@{}}
\toprule
Dataset & Ours\,/\,sig. & Sipers\,/\,sig. & GASF & GADF & XIRP \\
\midrule
WISDM & 0.0010 / 0.036 & 0.0043 / 0.083 & 0.0015 / 0.062 & 0.0013 / 0.066 & 0.0002 / 0.012 \\
PAMAP2 & 0.0015 / 0.044 & 0.0049 / 0.082 & 0.0016 / 0.065 & 0.0013 / 0.064 & 0.0002 / 0.011 \\
ESC-50 & 0.0017 / 0.042 & 0.0059 / 0.084 & 0.0016 / 0.062 & 0.0010 / 0.062 & 0.0002 / 0.011 \\
PTB-XL & 0.0016 / 0.051 & 0.0054 / 0.093 & 0.0030 / 0.074 & 0.0023 / 0.068 & 0.0002 / 0.011 \\
Swarm & 0.0126 / 0.062 & 0.0044 / 0.084 & 0.0026 / 0.063 & 0.0015 / 0.060 & 0.0002 / 0.012 \\
STEAD & 0.0012 / 0.032 & 0.0050 / 0.080 & 0.0014 / 0.061 & 0.0010 / 0.064 & 0.0002 / 0.011 \\
EMG & 0.0009 / 0.022 & 0.0040 / 0.075 & 0.0013 / 0.061 & 0.0022 / 0.063 & 0.0006 / 0.011 \\
\bottomrule
\end{tabular}

\end{table}

\section{Classification per dataset (RQ4)}\label{sec:s_classification}
Table~\ref{tab:s_cls_original} gives the score of the classifiers on the original windows at $T=128$, and Table~\ref{tab:s_cls_drop} the drop of each combination per dataset, averaged over classifiers and lengths; Table~\ref{M-tab:utility} of the main text aggregates them. Table~\ref{tab:s_esc10} gives the per-clip results on ESC-10.

\begin{table}[H]
\centering
\caption{Score on the original windows at $T=128$: accuracy, or macro-F1 for the multilabel PTB-XL.}\label{tab:s_cls_original}
\footnotesize
\begin{tabular}{@{}lrrrrrr@{}}
\toprule
Dataset & KNN & SVM & RF & LSTM & GRU & CNN \\
\midrule
WISDM & 0.467 & 0.722 & 0.750 & 0.652 & 0.709 & 0.746 \\
PAMAP2 & 0.540 & 0.651 & 0.689 & 0.541 & 0.596 & 0.589 \\
PTB-XL & 0.299 & 0.328 & 0.154 & 0.525 & 0.593 & 0.571 \\
STEAD & 0.511 & 0.490 & 0.658 & 0.517 & 0.511 & 0.583 \\
EMG & 0.166 & 0.409 & 0.770 & 0.652 & 0.675 & 0.813 \\
\bottomrule
\end{tabular}

\end{table}

\begin{table}[H]
\centering
\caption{Drop from the original score, averaged over the six classifiers and the three lengths.}\label{tab:s_cls_drop}
\scriptsize
\setlength{\tabcolsep}{2.5pt}
\resizebox{\textwidth}{!}{\begin{tabular}{@{}lrrrrrrrrrr@{}}
\toprule
Dataset & Hirata\,/\,uns. & Hirata\,/\,sig. & Sipers\,/\,uns. & Sipers\,/\,sig. & Ours\,/\,uns. & Ours\,/\,sig. & GASF & GADF & MTF & XIRP \\
\midrule
WISDM & 0.087 & 0.027 & 0.066 & 0.001 & 0.068 & 0.000 & 0.000 & 0.000 & 0.069 & 0.000 \\
PAMAP2 & 0.146 & 0.119 & 0.115 & 0.000 & 0.125 & 0.000 & -0.001 & 0.000 & 0.115 & 0.000 \\
PTB-XL & 0.221 & 0.152 & 0.215 & 0.000 & 0.210 & 0.000 & -0.001 & 0.000 & 0.223 & 0.000 \\
STEAD & 0.007 & 0.002 & 0.004 & 0.000 & 0.004 & 0.000 & 0.000 & 0.000 & 0.007 & 0.000 \\
EMG & 0.150 & 0.109 & 0.082 & 0.001 & 0.085 & 0.000 & 0.000 & 0.000 & 0.074 & 0.000 \\
\bottomrule
\end{tabular}
}
\end{table}

\begin{table}[H]
\centering
\caption{ESC-10: per-clip accuracy of the original, averaged over the six classifiers, and drop of each combination (3\,000 windows per block).}\label{tab:s_esc10}
\scriptsize
\setlength{\tabcolsep}{2pt}
\resizebox{\textwidth}{!}{\begin{tabular}{@{}llrrrrrrrrrr@{}}
\toprule
Block & Original & Hirata\,/\,uns. & Hirata\,/\,sig. & Sipers\,/\,uns. & Sipers\,/\,sig. & Ours\,/\,uns. & Ours\,/\,sig. & GASF & GADF & MTF & XIRP \\
\midrule
2\,kHz, $T=512$ & 0.312 & 0.081 & 0.069 & 0.028 & 0.001 & 0.026 & -0.001 & 0.000 & -0.001 & 0.125 & -0.001 \\
4\,kHz, $T=512$ & 0.339 & 0.076 & 0.063 & 0.023 & -0.003 & 0.025 & 0.000 & -0.002 & -0.002 & 0.118 & 0.000 \\
8\,kHz, $T=512$ & 0.311 & 0.061 & 0.043 & 0.017 & -0.001 & 0.019 & 0.001 & 0.001 & 0.000 & 0.066 & 0.000 \\
2\,kHz, $T=1024$ & 0.324 & 0.065 & 0.063 & 0.020 & 0.000 & 0.023 & 0.000 & 0.002 & 0.002 & 0.145 & 0.000 \\
4\,kHz, $T=1024$ & 0.306 & 0.077 & 0.072 & 0.017 & 0.003 & 0.026 & 0.003 & 0.000 & 0.001 & 0.106 & 0.000 \\
8\,kHz, $T=1024$ & 0.305 & 0.065 & 0.047 & 0.020 & 0.005 & 0.024 & 0.004 & 0.002 & 0.003 & 0.080 & 0.000 \\
\bottomrule
\end{tabular}
}
\end{table}

\section{Computational cost (RQ4)}\label{sec:s_cost}
Table~\ref{tab:s_cost} gives the single-thread cost of every combination against the window length. Times were measured on an otherwise idle Intel Core i7-14650HX (12 cores, 24 threads) under WSL2 (Linux 6.6), with Python 3.10.20, NumPy 1.26.4 and SciPy 1.14.0, one series at a time and multi-threaded linear algebra disabled; classifiers used scikit-learn 1.5.1 and TensorFlow 2.11.0.

\begin{table}[H]
\centering
\caption{Median single-thread time (ms) to encode, store and decode one channel (WISDM, 20 series per length), log--log slope of time against $T$, and peak traced memory at $T=1024$.}\label{tab:s_cost}
\footnotesize
\begin{tabular}{@{}lrrrrr@{}}
\toprule
Decoder & $T=64$ & $T=256$ & $T=1024$ & Exponent & MB at 1024 \\
\midrule
Hirata\,/\,uns. & 2.0 & 58.5 & 2871.2 & 2.64 & 87 \\
Hirata\,/\,sig. & 2.3 & 71.8 & 2880.1 & 2.59 & 87 \\
Sipers\,/\,uns. & 1.1 & 27.1 & 1172.7 & 2.47 & 112 \\
Sipers\,/\,sig. & 1.5 & 29.9 & 1193.1 & 2.45 & 112 \\
Ours\,/\,uns. & 0.6 & 13.1 & 398.2 & 2.35 & 72 \\
Ours\,/\,sig. & 1.3 & 25.1 & 552.0 & 2.22 & 72 \\
GASF & 0.1 & 1.0 & 12.0 & 1.85 & 32 \\
GADF & 0.1 & 1.0 & 13.0 & 1.72 & 32 \\
MTF & 0.5 & 1.5 & 12.8 & 1.22 & 24 \\
XIRP & 0.1 & 0.6 & 12.8 & 2.03 & 32 \\
\bottomrule
\end{tabular}

\end{table}

{\small
\begin{thebibliography}{99}
\setlength{\itemsep}{0pt}
\raggedright

\bibitem{misc_wisdm_smartphone_and_smartwatch_activity_and_biometrics_dataset__507}
Gary M. Weiss. ``WISDM Smartphone and Smartwatch Activity and Biometrics Dataset''. \textit{UCI Machine Learning Repository}, 2019. doi:\url{10.24432/C5HK59}.

\bibitem{reiss2012pamap2}
Attila Reiss and Didier Stricker. ``Introducing a new benchmarked dataset for activity monitoring''. In \textit{Proceedings of the 16th International Symposium on Wearable Computers (ISWC)}, pp. 108--109, 2012. doi:\url{10.1109/ISWC.2012.13}.

\bibitem{piczak2015esc}
Karol J. Piczak. ``ESC: Dataset for environmental sound classification''. In \textit{Proceedings of the 23rd ACM International Conference on Multimedia}, pp. 1015--1018, 2015. doi:\url{10.1145/2733373.2806390}.

\bibitem{wagner2020ptbxl}
Patrick Wagner, Nils Strodthoff, Ralf-Dieter Bousseljot, Dieter Kreiseler, Fatima I. Lunze, Wojciech Samek, and Tobias Schaeffter. ``PTB-XL, a large publicly available electrocardiography dataset''. \textit{Scientific Data}, vol. 7, article 154, 2020. doi:\url{10.1038/s41597-020-0495-6}.

\bibitem{friis2006swarm}
E. Friis-Christensen, H. Lühr, and G. Hulot. ``Swarm: A constellation to study the Earth's magnetic field''. \textit{Earth, Planets and Space}, vol. 58, no. 4, pp. 351--358, 2006. doi:\url{10.1186/BF03351933}.

\bibitem{mousavi2019stead}
S. Mostafa Mousavi, Yixiao Sheng, Weiqiang Zhu, and Gregory C. Beroza. ``STanford EArthquake Dataset (STEAD): A global data set of seismic signals for AI''. \textit{IEEE Access}, vol. 7, pp. 179464--179476, 2019. doi:\url{10.1109/ACCESS.2019.2947848}.

\bibitem{krilova2018emg}
N. Krilova, I. Kastalskiy, V. Kazantsev, V. A. Makarov, and S. Lobov. ``EMG Data for Gestures''. \textit{UCI Machine Learning Repository}, 2018. doi:\url{10.24432/C5ZP5C}.

\bibitem{Torgerson1952}
Warren S. Torgerson. ``Multidimensional scaling: I. Theory and method''. \textit{Psychometrika}, vol. 17, no. 4, pp. 401--419, 1952. doi:\url{10.1007/BF02288916}.

\bibitem{hirata2008reproduction}
Yoshito Hirata, Shunsuke Horai, and Kazuyuki Aihara. ``Reproduction of distance matrices and original time series from recurrence plots and their applications''. \textit{The European Physical Journal Special Topics}, vol. 164, no. 1, pp. 13--22, 2008. doi:\url{10.1140/epjst/e2008-00830-8}.

\bibitem{sipers2017robust}
Aloys Sipers, Paul J. A. Borm, and Ralf L. M. Peeters. ``Robust reconstruction of a signal from its unthresholded recurrence plot subject to disturbances''. \textit{Physics Letters A}, vol. 381, no. 6, pp. 604--615, 2017. doi:\url{10.1016/j.physleta.2016.12.028}.
\end{thebibliography}
}

\end{document}
